% series-figures-preamble.tex
%
% Shared requirements and styles for the Clebsch-series figure set.
% A host manuscript inputs this once in its preamble, then inputs any of
% the individual figure files in its body.
%
% The styles are deliberately colour-free in meaning: every distinction is
% carried by line style, line weight, or arrowhead.  The one tint used
% (light grey) marks the "you are here" node of the series map and remains
% readable in black and white.

\usepackage{amsmath,amssymb}
\usepackage{array}
\usepackage{booktabs}
\usepackage{tikz}
\usetikzlibrary{arrows.meta,calc,positioning,fit,backgrounds,
  decorations.pathreplacing}

\tikzset{
  % Ordinary labelled box.
  csfnode/.style={draw,thin,rounded corners=2pt,align=center,
    inner sep=4pt,font=\small},
  % Emphasised box: heavier rule plus a light tint.
  csfhi/.style={line width=1pt,fill=black!8},
  % The two states of a "you are here" node.  These exist as single style
  % names because a key list does not re-split an expanded macro.
  csfplain/.style={csfnode},
  csfmarked/.style={csfnode,csfhi},
  % The central object of a diagram.
  csfcentre/.style={csfnode,line width=0.8pt},
  % Standard directed connector.
  csfarr/.style={-{Latex[length=2mm]},semithick},
  % Bidirectional connector (a passage that forgets outward and recovers
  % inward, or an identity read in both directions).
  csfbiarr/.style={{Latex[length=2mm]}-{Latex[length=2mm]},semithick},
  % An asserted-but-unproved bridge.
  csfopen/.style={-{Latex[length=2mm]},semithick,dashed},
  % A proved statement that is not yet in the manuscript.  Deliberately a
  % DOUBLE rule: it must never be confused with the dashed csfopen bridge,
  % which asserts something no theorem supports.  This one is established
  % mathematics awaiting promotion, so it is drawn heavier than an ordinary
  % arrow, not lighter.
  csfforth/.style={-{Implies},double,double distance=1.6pt,line width=0.5pt},
  % The tag that names that status, ruled to match csfforth.
  csfbadge/.style={draw,double,double distance=1pt,line width=0.4pt,
    rounded corners=1pt,align=center,font=\scriptsize\bfseries,
    inner sep=3pt,fill=white},
  % A reformulation rather than a new identity.
  csfweak/.style={thin,densely dotted},
  % The one link in a diagram carrying real content.
  csfstrong/.style={{Latex[length=2mm]}-{Latex[length=2mm]},line width=1pt},
  % Edge label.
  csflab/.style={font=\scriptsize,fill=white,inner sep=1.2pt,align=center},
  csflabsmall/.style={font=\tiny,fill=white,inner sep=1pt,align=center},
}

% ---------------------------------------------------------------------------
% "You are here" selection used by the series map (D1).
%
%   \csfSetHighlight{III}   % or I, II, IV, V, or none
%
% After this call \csfHL{I} expands to the TikZ style name for Paper I's
% node: either "csfplain" or "csfmarked".  The series map macro calls
% \csfSetHighlight itself, so a document normally never calls it directly.
% ---------------------------------------------------------------------------
\providecommand{\csfHL}[1]{\csname csfhl@#1\endcsname}
\providecommand{\csfSetHighlight}[1]{%
  \expandafter\gdef\csname csfhl@I\endcsname{csfplain}%
  \expandafter\gdef\csname csfhl@II\endcsname{csfplain}%
  \expandafter\gdef\csname csfhl@III\endcsname{csfplain}%
  \expandafter\gdef\csname csfhl@IV\endcsname{csfplain}%
  \expandafter\gdef\csname csfhl@V\endcsname{csfplain}%
  \expandafter\gdef\csname csfhl@none\endcsname{csfplain}%
  \expandafter\gdef\csname csfhl@#1\endcsname{csfmarked}%
}
\csfSetHighlight{none}

% ---------------------------------------------------------------------------
% Return-arrow selection used by the causal spine (D6).
%
%   \csfSetSpineReturn{current}   % the manuscript's reconstruction theorem
%   \csfSetSpineReturn{final}     % the triangle--Pfaffian recognition
%
% After this call \csfSpineFinal expands to 1 for the "final" variant and to
% 0 otherwise, so an unrecognised argument falls back to "current".  The
% spine macro calls \csfSetSpineReturn itself.
% ---------------------------------------------------------------------------
\providecommand{\csfSpineFinal}{\csname csfspine@final\endcsname}
\providecommand{\csfSetSpineReturn}[1]{%
  \expandafter\gdef\csname csfspine@current\endcsname{0}%
  \expandafter\gdef\csname csfspine@final\endcsname{0}%
  \expandafter\gdef\csname csfspine@#1\endcsname{1}%
}
\csfSetSpineReturn{current}

\providecommand{\PGtwo}[1]{\operatorname{PG}(2,#1)}
